\documentclass[11pt]{article}
\usepackage{amsmath,amssymb,amsthm,mathtools,enumitem,booktabs}
\usepackage[margin=1in]{geometry}

\newcommand{\Psym}{\mathsf P_\infty}
\newcommand{\Mz}{M_\zeta}
\newcommand{\Uinf}{\mathcal U_\infty}
\newcommand{\Blegal}{\mathcal B}

\title{Step 174: Evaluation Attempt for \(L_{\rho,k}(G)\)}
\author{Six Birds Foundations III}
\date{2026-05-15}

\begin{document}
\maketitle

\section{Verdict}

\[
\boxed{\texttt{L\_verdict = V\_L\_symbolic\_form}.}
\]

This step evaluates \(L_{\rho_1,0}(G_*)\) for a concrete nonzero legal Burnol
generator to an explicit PSWF-series formula.  It does not prove that the value
is zero or nonzero.  In this evaluation, zero/nonzero is undecided.

\section{Inherited Evaluator Convention}

Step 158 declares the pulled-evaluator context
\[
H_\eta=\overline{\operatorname{span}\{\eta^a_{\rho,k}:
\eta^a_{\rho,k}=J_a^*Y^a_{\rho,k}\}}.
\]
Step 162 records \(P_\eta\) as the projection onto this closed pulled
zero-evaluator span.  Step 162 does not introduce a new normalization for
\(y_{\rho,k}\).

Step 153 gives the actual evaluator convention:
\[
[f,Y^a_{w,k}]_{L_a}=M(f)^{(k)}(w),
\qquad f\in L_a.
\]
In the Hilbert convention,
\[
\mathcal M_\Gamma(Y^a_{w,k})(s)
=\partial_{\overline w}^{\,k}K_a^\Gamma(s,w),
\]
and under Burnol's bilinear convention the derivative is \(\partial_w^k\).
For this step, at \(k=0\), the convention issue is only conjugation placement;
we use the Hilbert convention
\[
\langle f,y_{\rho,0}\rangle=f(\rho).
\]
Thus
\[
L_{\rho,0}(G)=(\Psym\Mz G)(\rho).
\]

\section{Concrete Legal Generator}

Set
\[
a=1,\qquad A=4,\qquad \lambda=1,
\qquad
\rho_1=\frac12+14.134725141734693\,i.
\]
The inherited records do not fix \(\lambda\), so \(\lambda=1\) is a concrete
normalization choice for this data point.

Let
\[
\beta(u)=
\begin{cases}
e^{-1/(1-u^2)},& |u|<1,\\
0,& |u|\ge1.
\end{cases}
\]
With \(\varepsilon=1/5\) and centers
\[
c_1=\frac32,\qquad c_2=\frac52,\qquad c_3=\frac72,
\]
define
\[
b_j(t)=\beta\left(\frac{t-c_j}{\varepsilon}\right).
\]
These are nonzero \(C_c^\infty(1,4)\) bumps with disjoint supports.  Let
\[
A_j=\int_1^4 b_j(t)\,dt,\qquad
B_j=\int_1^4 b_j(t)t^{-1}\,dt,
\]
and
\[
D=B_2A_3-A_2B_3.
\]
Since the supports have different \(t^{-1}\)-averages, \(D\ne0\).  Set
\[
\alpha=\frac{A_3B_1-A_1B_3}{D},
\qquad
\beta_3=\frac{A_2B_1-A_1B_2}{D},
\]
and
\[
g_*(t)=b_1(t)-\alpha b_2(t)+\beta_3 b_3(t).
\]
Then
\[
\int_1^4 g_*(t)\,dt=0,\qquad
\int_1^4 g_*(t)t^{-1}\,dt=0.
\]
Therefore the right Mellin transform
\[
G_*(s)=\widehat g_*(s)=\int_1^4 g_*(t)t^{-s}\,dt
\]
satisfies the Step 119 endpoint vanishings
\[
G_*(0)=0,\qquad G_*(1)=0.
\]
It is a nontrivial legal Burnol generator.

\section{T3a. Zeta-Multiplied Profile}

Define
\[
F_*(s)=(\Mz G_*)(s)=\zeta(s)G_*(s).
\]
At the chosen zero,
\[
F_*(\rho_1)=\zeta(\rho_1)G_*(\rho_1)=0.
\]

\section{T3b. Apply the Step 173 Operational Identity}

Step 173 gives, for \(s=1/2+i\tau\),
\[
(\Psym F)(s)
=
F(s)
-
\int_{\mathbb R}
\frac{\sin(\tau-u)}{\pi(\tau-u)}
F(1/2+iu)\,du
-
\sum_{n\ge0}\Psi_n(s)\langle F,\Psi_n\rangle,
\]
where \(\lambda=1\) and
\[
\Psi_n(s)=
\Uinf\left(
\frac{(I-P_1)\phi_n^1}{\sqrt{1-\mu_n^1}}
\right)(s).
\]
Here \(\phi_n^1\) are the Slepian--Pollak PSWFs diagonalizing
\[
\widehat P_1P_1\widehat P_1\phi_n^1=\mu_n^1\phi_n^1.
\]
Applying this to \(F_*\) gives the three-term decomposition:
\[
(\Psym F_*)(s)=D_*(s)-C_*(s)-R_*(s),
\]
with
\[
D_*(s)=F_*(s),
\]
\[
C_*(1/2+i\tau)=
\int_{\mathbb R}
\frac{\sin(\tau-u)}{\pi(\tau-u)}
\zeta(1/2+iu)G_*(1/2+iu)\,du,
\]
and
\[
R_*(s)=
\sum_{n\ge0}\Psi_n(s)
\int_{\mathbb R}
\zeta(1/2+iu)G_*(1/2+iu)
\overline{\Psi_n(1/2+iu)}\,du.
\]
The displayed leading truncation through \(n=2\) is
\[
R_*^{[0,2]}(s)=
\sum_{n=0}^{2}\Psi_n(s)
\int_{\mathbb R}
\zeta(1/2+iu)G_*(1/2+iu)
\overline{\Psi_n(1/2+iu)}\,du.
\]
The remainder is
\[
R_*^{[>2]}(s)=R_*(s)-R_*^{[0,2]}(s).
\]

\section{T3c. Pair with the Zero Evaluator}

For \(k=0\),
\[
L_{\rho_1,0}(G_*)=(\Psym F_*)(\rho_1).
\]
Since \(D_*(\rho_1)=0\), the symbolic value is
\[
\boxed{
L_{\rho_1,0}(G_*)
=
-C_*(\rho_1)-R_*(\rho_1).
}
\]
Writing \(\gamma_1=14.134725141734693\), this is
\[
\boxed{
\begin{aligned}
L_{\rho_1,0}(G_*)
&=
-\int_{\mathbb R}
\frac{\sin(\gamma_1-u)}{\pi(\gamma_1-u)}
\zeta(1/2+iu)G_*(1/2+iu)\,du\\
&\quad
-\sum_{n\ge0}\Psi_n(\rho_1)
\int_{\mathbb R}
\zeta(1/2+iu)G_*(1/2+iu)
\overline{\Psi_n(1/2+iu)}\,du.
\end{aligned}
}
\]
The first three PSWF terms are obtained by replacing the infinite sum by
\(\sum_{n=0}^{2}\), with remainder \(R_*^{[>2]}(\rho_1)\).

\section{T3d. Parameter Choices}

The concrete choices used are:
\[
\lambda=1,\qquad \rho=\rho_1,\qquad k=0,\qquad
\text{PSWF leading block } n=0,1,2.
\]
No inherited record fixes \(\lambda\); \(\lambda=1\) is a declared data-point
normalization.

\section{Branch C Implication}

The result is a symbolic formula, not an annihilation proof and not an
obstruction proof.  Branch C is advanced from an undefined matrix element to a
specific PSWF/sinc integral expression for one nontrivial legal generator.
Numerical sign or nonvanishing now depends on evaluating the PSWF transforms
\(\Psi_n\) and the zeta-weighted integrals.

\section{Nonclaim Boundary}

This step does not prove RH, does not prove \(\Xi_{\mathrm{BC}}=0\), does not
prove full ZI-COV, does not prove projected jet-surjectivity, and does not
decide whether \(L_{\rho_1,0}(G_*)\) is zero or nonzero.  It does not weaken
any retained no-go and does not replace the inherited Sonin/prolate projection
by a zeta-zero spectral projection.

\end{document}
